\IfFileExists{stacks-project.cls}{%
\documentclass{stacks-project}
}{%
\documentclass{amsart}
}

% For dealing with references we use the comment environment
\usepackage{verbatim}
\newenvironment{reference}{\comment}{\endcomment}
%\newenvironment{reference}{}{}
\newenvironment{slogan}{\comment}{\endcomment}
\newenvironment{history}{\comment}{\endcomment}

% For commutative diagrams we use Xy-pic
\usepackage[all]{xy}

% We use 2cell for 2-commutative diagrams.
\xyoption{2cell}
\UseAllTwocells

% We use multicol for the list of chapters between chapters
\usepackage{multicol}

% This is generally recommended for better output
\usepackage{lmodern}
\usepackage[T1]{fontenc}

% For cross-file-references
\usepackage{xr-hyper}

% Package for hypertext links:
\usepackage{hyperref}

% For any local file, say "hello.tex" you want to link to please
% use \externaldocument[hello-]{hello}
\externaldocument[introduction-]{introduction}
\externaldocument[conventions-]{conventions}
\externaldocument[sets-]{sets}
\externaldocument[categories-]{categories}
\externaldocument[topology-]{topology}
\externaldocument[sheaves-]{sheaves}
\externaldocument[sites-]{sites}
\externaldocument[stacks-]{stacks}
\externaldocument[fields-]{fields}
\externaldocument[algebra-]{algebra}
\externaldocument[brauer-]{brauer}
\externaldocument[homology-]{homology}
\externaldocument[derived-]{derived}
\externaldocument[simplicial-]{simplicial}
\externaldocument[more-algebra-]{more-algebra}
\externaldocument[smoothing-]{smoothing}
\externaldocument[modules-]{modules}
\externaldocument[sites-modules-]{sites-modules}
\externaldocument[injectives-]{injectives}
\externaldocument[cohomology-]{cohomology}
\externaldocument[sites-cohomology-]{sites-cohomology}
\externaldocument[dga-]{dga}
\externaldocument[dpa-]{dpa}
\externaldocument[sdga-]{sdga}
\externaldocument[hypercovering-]{hypercovering}
\externaldocument[schemes-]{schemes}
\externaldocument[constructions-]{constructions}
\externaldocument[properties-]{properties}
\externaldocument[morphisms-]{morphisms}
\externaldocument[coherent-]{coherent}
\externaldocument[divisors-]{divisors}
\externaldocument[limits-]{limits}
\externaldocument[varieties-]{varieties}
\externaldocument[topologies-]{topologies}
\externaldocument[descent-]{descent}
\externaldocument[perfect-]{perfect}
\externaldocument[more-morphisms-]{more-morphisms}
\externaldocument[flat-]{flat}
\externaldocument[groupoids-]{groupoids}
\externaldocument[more-groupoids-]{more-groupoids}
\externaldocument[etale-]{etale}
\externaldocument[chow-]{chow}
\externaldocument[intersection-]{intersection}
\externaldocument[pic-]{pic}
\externaldocument[weil-]{weil}
\externaldocument[adequate-]{adequate}
\externaldocument[dualizing-]{dualizing}
\externaldocument[duality-]{duality}
\externaldocument[discriminant-]{discriminant}
\externaldocument[derham-]{derham}
\externaldocument[local-cohomology-]{local-cohomology}
\externaldocument[algebraization-]{algebraization}
\externaldocument[curves-]{curves}
\externaldocument[resolve-]{resolve}
\externaldocument[models-]{models}
\externaldocument[functors-]{functors}
\externaldocument[equiv-]{equiv}
\externaldocument[pione-]{pione}
\externaldocument[etale-cohomology-]{etale-cohomology}
\externaldocument[proetale-]{proetale}
\externaldocument[relative-cycles-]{relative-cycles}
\externaldocument[more-etale-]{more-etale}
\externaldocument[trace-]{trace}
\externaldocument[crystalline-]{crystalline}
\externaldocument[spaces-]{spaces}
\externaldocument[spaces-properties-]{spaces-properties}
\externaldocument[spaces-morphisms-]{spaces-morphisms}
\externaldocument[decent-spaces-]{decent-spaces}
\externaldocument[spaces-cohomology-]{spaces-cohomology}
\externaldocument[spaces-limits-]{spaces-limits}
\externaldocument[spaces-divisors-]{spaces-divisors}
\externaldocument[spaces-over-fields-]{spaces-over-fields}
\externaldocument[spaces-topologies-]{spaces-topologies}
\externaldocument[spaces-descent-]{spaces-descent}
\externaldocument[spaces-perfect-]{spaces-perfect}
\externaldocument[spaces-more-morphisms-]{spaces-more-morphisms}
\externaldocument[spaces-flat-]{spaces-flat}
\externaldocument[spaces-groupoids-]{spaces-groupoids}
\externaldocument[spaces-more-groupoids-]{spaces-more-groupoids}
\externaldocument[bootstrap-]{bootstrap}
\externaldocument[spaces-pushouts-]{spaces-pushouts}
\externaldocument[spaces-chow-]{spaces-chow}
\externaldocument[groupoids-quotients-]{groupoids-quotients}
\externaldocument[spaces-more-cohomology-]{spaces-more-cohomology}
\externaldocument[spaces-simplicial-]{spaces-simplicial}
\externaldocument[spaces-duality-]{spaces-duality}
\externaldocument[formal-spaces-]{formal-spaces}
\externaldocument[restricted-]{restricted}
\externaldocument[spaces-resolve-]{spaces-resolve}
\externaldocument[formal-defos-]{formal-defos}
\externaldocument[defos-]{defos}
\externaldocument[cotangent-]{cotangent}
\externaldocument[examples-defos-]{examples-defos}
\externaldocument[algebraic-]{algebraic}
\externaldocument[examples-stacks-]{examples-stacks}
\externaldocument[stacks-sheaves-]{stacks-sheaves}
\externaldocument[criteria-]{criteria}
\externaldocument[artin-]{artin}
\externaldocument[quot-]{quot}
\externaldocument[stacks-properties-]{stacks-properties}
\externaldocument[stacks-morphisms-]{stacks-morphisms}
\externaldocument[stacks-limits-]{stacks-limits}
\externaldocument[stacks-cohomology-]{stacks-cohomology}
\externaldocument[stacks-perfect-]{stacks-perfect}
\externaldocument[stacks-introduction-]{stacks-introduction}
\externaldocument[stacks-more-morphisms-]{stacks-more-morphisms}
\externaldocument[stacks-geometry-]{stacks-geometry}
\externaldocument[moduli-]{moduli}
\externaldocument[moduli-curves-]{moduli-curves}
\externaldocument[examples-]{examples}
\externaldocument[exercises-]{exercises}
\externaldocument[guide-]{guide}
\externaldocument[desirables-]{desirables}
\externaldocument[coding-]{coding}
\externaldocument[obsolete-]{obsolete}
\externaldocument[fdl-]{fdl}
\externaldocument[index-]{index}

% Theorem environments.
%
\theoremstyle{plain}
\newtheorem{theorem}[subsection]{Theorem}
\newtheorem{proposition}[subsection]{Proposition}
\newtheorem{lemma}[subsection]{Lemma}

\theoremstyle{definition}
\newtheorem{definition}[subsection]{Definition}
\newtheorem{example}[subsection]{Example}
\newtheorem{exercise}[subsection]{Exercise}
\newtheorem{situation}[subsection]{Situation}

\theoremstyle{remark}
\newtheorem{remark}[subsection]{Remark}
\newtheorem{remarks}[subsection]{Remarks}

\numberwithin{equation}{subsection}

% Macros
%
\def\lim{\mathop{\mathrm{lim}}\nolimits}
\def\colim{\mathop{\mathrm{colim}}\nolimits}
\def\Spec{\mathop{\mathrm{Spec}}}
\def\Hom{\mathop{\mathrm{Hom}}\nolimits}
\def\Ext{\mathop{\mathrm{Ext}}\nolimits}
\def\SheafHom{\mathop{\mathcal{H}\!\mathit{om}}\nolimits}
\def\SheafExt{\mathop{\mathcal{E}\!\mathit{xt}}\nolimits}
\def\Sch{\mathit{Sch}}
\def\Mor{\mathop{\mathrm{Mor}}\nolimits}
\def\Ob{\mathop{\mathrm{Ob}}\nolimits}
\def\Sh{\mathop{\mathit{Sh}}\nolimits}
\def\NL{\mathop{N\!L}\nolimits}
\def\CH{\mathop{\mathrm{CH}}\nolimits}
\def\proetale{{pro\text{-}\acute{e}tale}}
\def\etale{{\acute{e}tale}}
\def\QCoh{\mathit{QCoh}}
\def\Ker{\mathop{\mathrm{Ker}}}
\def\Im{\mathop{\mathrm{Im}}}
\def\Coker{\mathop{\mathrm{Coker}}}
\def\Coim{\mathop{\mathrm{Coim}}}

% Boxtimes
%
\DeclareMathSymbol{\boxtimes}{\mathbin}{AMSa}{"02}

%
% Macros for moduli stacks/spaces
%
\def\QCohstack{\mathcal{QC}\!\mathit{oh}}
\def\Cohstack{\mathcal{C}\!\mathit{oh}}
\def\Spacesstack{\mathcal{S}\!\mathit{paces}}
\def\Quotfunctor{\mathrm{Quot}}
\def\Hilbfunctor{\mathrm{Hilb}}
\def\Curvesstack{\mathcal{C}\!\mathit{urves}}
\def\Polarizedstack{\mathcal{P}\!\mathit{olarized}}
\def\Complexesstack{\mathcal{C}\!\mathit{omplexes}}
% \Pic is the operator that assigns to X its picard group, usage \Pic(X)
% \Picardstack_{X/B} denotes the Picard stack of X over B
% \Picardfunctor_{X/B} denotes the Picard functor of X over B
\def\Pic{\mathop{\mathrm{Pic}}\nolimits}
\def\Picardstack{\mathcal{P}\!\mathit{ic}}
\def\Picardfunctor{\mathrm{Pic}}
\def\Deformationcategory{\mathcal{D}\!\mathit{ef}}
\usepackage{cleveref}
\crefname{originallemma}{original lemma}{original lemmas}
\Crefname{originallemma}{Original lemma}{Original lemmas}
\makeatletter
\newlabel{lemma-directed-inverse-system-has-limit}{{07SF}{source}{}{sourcepointer-1}{}}
\newlabel{lemma-directed-inverse-system-has-limit@cref}{{[originallemma][0][]07SF}{[1][0][]source}{}{}{}}
\makeatother
\crefname{originalreference}{original reference}{original references}
\Crefname{originalreference}{Original reference}{Original references}
\makeatletter
\newlabel{limits-lemma-descend-etale}{{07SL}{source}{}{sourcepointer-2}{}}
\newlabel{limits-lemma-descend-etale@cref}{{[originalreference][0][]07SL}{[1][0][]source}{}{}{}}
\makeatother
\crefname{originalreference}{original reference}{original references}
\Crefname{originalreference}{Original reference}{Original references}
\makeatletter
\newlabel{limits-lemma-descend-finite-presentation}{{07SK}{source}{}{sourcepointer-3}{}}
\newlabel{limits-lemma-descend-finite-presentation@cref}{{[originalreference][0][]07SK}{[1][0][]source}{}{}{}}
\makeatother
\crefname{originalreference}{original reference}{original references}
\Crefname{originalreference}{Original reference}{Original references}
\makeatletter
\newlabel{limits-lemma-descend-monomorphism}{{07SM}{source}{}{sourcepointer-4}{}}
\newlabel{limits-lemma-descend-monomorphism@cref}{{[originalreference][0][]07SM}{[1][0][]source}{}{}{}}
\makeatother
\crefname{originalproposition}{original proposition}{original propositions}
\Crefname{originalproposition}{Original proposition}{Original propositions}
\makeatletter
\newlabel{proposition-characterize-locally-finite-presentation}{{04AK}{source}{}{sourcepointer-5}{}}
\newlabel{proposition-characterize-locally-finite-presentation@cref}{{[originalproposition][0][]04AK}{[1][0][]source}{}{}{}}
\makeatother
\crefname{originalreference}{original reference}{original references}
\Crefname{originalreference}{Original reference}{Original references}
\makeatletter
\newlabel{spaces-properties-lemma-quasi-compact-affine-cover}{{03H6}{source}{}{sourcepointer-6}{}}
\newlabel{spaces-properties-lemma-quasi-compact-affine-cover@cref}{{[originalreference][0][]03H6}{[1][0][]source}{}{}{}}
\makeatother
\crefname{originalreference}{original reference}{original references}
\Crefname{originalreference}{Original reference}{Original references}
\makeatletter
\newlabel{spaces-section-presentations}{{0261}{source}{}{sourcepointer-7}{}}
\newlabel{spaces-section-presentations@cref}{{[originalreference][0][]0261}{[1][0][]source}{}{}{}}
\makeatother
\crefname{originalreference}{original reference}{original references}
\Crefname{originalreference}{Original reference}{Original references}
\makeatletter
\newlabel{spaces-theorem-presentation}{{02WW}{source}{}{sourcepointer-8}{}}
\newlabel{spaces-theorem-presentation@cref}{{[originalreference][0][]02WW}{[1][0][]source}{}{}{}}
\makeatother
\title{Finite presentation spaces over a limit form the colimit category}
\author{The Stacks Project Authors}
\date{stacks-project:spaces-limits.tex:local-snapshot-20261012; GNU FDL 1.2 or later}
\begin{document}
\maketitle
\noindent\textit{Editorial note.} Counted span is the author's own proof environment, source lines 2161-2227, following the statement at 2147-2159.\par\medskip
\noindent\textit{Editorial note.} The single '(details omitted)' covers finite presentation of W over U and of R' over R: etale maps are locally of finite presentation and fp is preserved under base change; flagged in the notes.\par\medskip
\noindent\textit{Editorial note.} Steps the author leaves to the reader are flagged here for the reviewer: source line(s) 2179 carry an "omitted"/"follows easily"/"immediate from the definitions" phrase. Those sentences are retained verbatim; the displayed argument around them stays inside the counted span.\par\medskip
\noindent\textit{Allowed standard prerequisites (editorial).}\begin{itemize}
\item Affine etale cover of a quasi-compact algebraic space. Properties of Spaces, Lemma spaces-properties-lemma-quasi-compact-affine-cover (cross-chapter, cited at \textasciitilde{}2163-2164). Produces the affine etale U\_0 -> X\_0.
\item Presentation of the limit space. lemma-directed-inverse-system-has-limit (same chapter, cited at \textasciitilde{}2168-2169). Gives X = U/R with U = lim U\_i and R = lim R\_i.
\item Descent of finite presentation along affine limits of schemes. Limits, Lemma limits-lemma-descend-finite-presentation (cross-chapter, cited at \textasciitilde{}2180 and \textasciitilde{}2190). Descends W -> U and R' -> R to some finite level; applied twice.
\item Descent of etale maps and monomorphisms. Limits, Lemmas limits-lemma-descend-etale and limits-lemma-descend-monomorphism (cross-chapter, cited at \textasciitilde{}2193-2194). Makes s'\_i, t'\_i etale and j'\_i a monomorphism after increasing i.
\item Presentation theorem for algebraic spaces. Spaces, Section spaces-section-presentations and Theorem spaces-theorem-presentation (cross-chapter, cited at \textasciitilde{}2177 and \textasciitilde{}2199). Y = W/R' and Y\_i = W\_i/R'\_i are algebraic spaces.
\end{itemize}\medskip
\ifdefined\hypertarget\hypertarget{sourcepointer-1}{}\fi
\noindent\textit{Source reference.} Original lemma 07SF: sources/stacks-spaces-limits/pinned-original-source.tex, source lines 878--878. Navigation pointer only: the counted proof cites this statement by \textbackslash{}ref as a standing input; the deduction the counted proof has to make is carried out inside the counted span, and no part of that deduction is delegated to this pointer.\par\medskip
\ifdefined\hypertarget\hypertarget{sourcepointer-2}{}\fi
\noindent\textit{Source reference.} Original reference 07SL: focused-seeds/geometry/S9-spaces-limits-descend-finite-presentation/ref-tags.txt, source lines 5--5. Navigation pointer only: the counted proof cites this statement by \textbackslash{}ref as a standing input; the deduction the counted proof has to make is carried out inside the counted span, and no part of that deduction is delegated to this pointer.\par\medskip
\ifdefined\hypertarget\hypertarget{sourcepointer-3}{}\fi
\noindent\textit{Source reference.} Original reference 07SK: focused-seeds/geometry/S9-spaces-limits-descend-finite-presentation/ref-tags.txt, source lines 6--6. Navigation pointer only: the counted proof cites this statement by \textbackslash{}ref as a standing input; the deduction the counted proof has to make is carried out inside the counted span, and no part of that deduction is delegated to this pointer.\par\medskip
\ifdefined\hypertarget\hypertarget{sourcepointer-4}{}\fi
\noindent\textit{Source reference.} Original reference 07SM: focused-seeds/geometry/S9-spaces-limits-descend-finite-presentation/ref-tags.txt, source lines 7--7. Navigation pointer only: the counted proof cites this statement by \textbackslash{}ref as a standing input; the deduction the counted proof has to make is carried out inside the counted span, and no part of that deduction is delegated to this pointer.\par\medskip
\ifdefined\hypertarget\hypertarget{sourcepointer-5}{}\fi
\noindent\textit{Source reference.} Original proposition 04AK: sources/stacks-spaces-limits/pinned-original-source.tex, source lines 694--694. Navigation pointer only: the counted proof cites this statement by \textbackslash{}ref as a standing input; the deduction the counted proof has to make is carried out inside the counted span, and no part of that deduction is delegated to this pointer.\par\medskip
\ifdefined\hypertarget\hypertarget{sourcepointer-6}{}\fi
\noindent\textit{Source reference.} Original reference 03H6: focused-seeds/geometry/S9-spaces-limits-descend-finite-presentation/ref-tags.txt, source lines 8--8. Navigation pointer only: the counted proof cites this statement by \textbackslash{}ref as a standing input; the deduction the counted proof has to make is carried out inside the counted span, and no part of that deduction is delegated to this pointer.\par\medskip
\ifdefined\hypertarget\hypertarget{sourcepointer-7}{}\fi
\noindent\textit{Source reference.} Original reference 0261: focused-seeds/geometry/S9-spaces-limits-descend-finite-presentation/ref-tags.txt, source lines 9--9. Navigation pointer only: the counted proof cites this statement by \textbackslash{}ref as a standing input; the deduction the counted proof has to make is carried out inside the counted span, and no part of that deduction is delegated to this pointer.\par\medskip
\ifdefined\hypertarget\hypertarget{sourcepointer-8}{}\fi
\noindent\textit{Source reference.} Original reference 02WW: focused-seeds/geometry/S9-spaces-limits-descend-finite-presentation/ref-tags.txt, source lines 10--10. Navigation pointer only: the counted proof cites this statement by \textbackslash{}ref as a standing input; the deduction the counted proof has to make is carried out inside the counted span, and no part of that deduction is delegated to this pointer.\par\medskip
\begin{lemma}
\label{lemma-descend-finite-presentation}
Let $S$ be a scheme. Let $I$ be a directed set.
Let $(X_i, f_{ii'})$ be an inverse system over $I$ of algebraic spaces
over $S$. Assume
\begin{enumerate}
\item the morphisms $f_{ii'} : X_i \to X_{i'}$ are affine,
\item the spaces $X_i$ are quasi-compact and quasi-separated.
\end{enumerate}
Let $X = \lim_i X_i$. Then the category of algebraic spaces
of finite presentation over $X$ is the colimit over $I$ of the
categories of algebraic spaces of finite presentation over $X_i$.
\end{lemma}

\begin{proof}
Pick $0 \in I$. Choose a surjective \'etale morphism $U_0 \to X_0$ where
$U_0$ is an affine scheme (Properties of Spaces, Lemma
\ref{spaces-properties-lemma-quasi-compact-affine-cover}).
Set $U_i = X_i \times_{X_0} U_0$. Set $R_0 = U_0 \times_{X_0} U_0$ and
$R_i = R_0 \times_{X_0} X_i$. Denote $s_i, t_i : R_i \to U_i$ and
$s, t : R \to U$ the two projections. In the proof of
Lemma \ref{lemma-directed-inverse-system-has-limit} we have
seen that there exists a presentation $X = U/R$ with
$U = \lim U_i$ and $R = \lim R_i$. Note that $U_i$ and $U$ are affine and
that $R_i$ and $R$ are quasi-compact and separated (as $X_i$ is
quasi-separated). Let $Y$ be an algebraic space over $S$ and let
$Y \to X$ be a morphism of finite presentation. Set $V = U \times_X Y$.
This is an algebraic space of finite presentation over $U$.
Choose an affine scheme $W$ and a surjective \'etale morphism $W \to V$.
Then $W \to Y$ is surjective \'etale as well. Set $R' = W \times_Y W$
so that $Y = W/R'$ (see Spaces, Section \ref{spaces-section-presentations}).
Note that $W$ is a scheme of finite presentation over $U$ and that $R'$
is a scheme of finite presentation over $R$ (details omitted).
By Limits, Lemma \ref{limits-lemma-descend-finite-presentation}
we can find an index $i$ and a morphism of schemes $W_i \to U_i$ of
finite presentation whose base change to $U$ gives $W \to U$. Similarly
we can find, after possibly increasing $i$, a scheme $R'_i$ of finite
presentation over $R_i$ whose base change to $R$ is $R'$.
The projection morphisms $s', t' : R' \to W$ are morphisms over
the projection morphisms $s, t : R \to U$. Hence we can view $s'$,
resp.\ $t'$ as a morphism between schemes of finite presentation over
$U$ (with structure morphism $R' \to U$ given by $R' \to R$ followed
by $s$, resp.\ $t$). Hence we can apply
Limits, Lemma \ref{limits-lemma-descend-finite-presentation}
again to see that, after possibly increasing $i$, there exist
morphisms $s'_i, t'_i : R'_i \to W_i$, whose base change to $U$
is $S', t'$. By Limits, Lemmas \ref{limits-lemma-descend-etale} and
\ref{limits-lemma-descend-monomorphism}
we may assume that $s'_i, t'_i$ are \'etale and that
$j'_i : R'_i \to W_i \times_{X_i} W_i$ is a monomorphism (here we
view $j'_i$ as a morphism of schemes of finite presentation over $U_i$ via
one of the projections -- it doesn't matter which one). Setting
$Y_i = W_i/R'_i$ (see Spaces, Theorem \ref{spaces-theorem-presentation})
we obtain an algebraic space of finite presentation
over $X_i$ whose base change to $X$ is isomorphic to $Y$.

\medskip\noindent
This shows that every algebraic space of finite presentation over $X$ comes
from an algebraic space of finite presentation over some $X_i$, i.e.,
it shows that the functor of the lemma is essentially surjective. To
show that it is fully faithful, consider an index $0 \in I$ and two
algebraic spaces $Y_0, Z_0$ of finite presentation over $X_0$.
Set $Y_i = X_i \times_{X_0} Y_0$, $Y = X \times_{X_0} Y_0$,
$Z_i = X_i \times_{X_0} Z_0$, and $Z = X \times_{X_0} Z_0$. Let
$\alpha : Y \to Z$ be a morphism of algebraic spaces over $X$.
Choose a surjective \'etale morphism $V_0 \to Y_0$ where $V_0$ is
an affine scheme. Set $V_i = V_0 \times_{Y_0} Y_i$ and
$V = V_0 \times_{Y_0} Y$ which are affine schemes endowed with
surjective \'etale morphisms to $Y_i$ and $Y$. The composition
$V \to Y \to Z \to Z_0$ comes from a (essentially unique) morphism
$V_i \to Z_0$ for some $i \geq 0$ by
Proposition \ref{proposition-characterize-locally-finite-presentation}
(applied to $Z_0 \to X_0$ which is of finite presentation by assumption).
After increasing $i$ the two compositions
$$
V_i \times_{Y_i} V_i \to V_i \to Z_0
$$
are equal as this is true in the limit. Hence we obtain a (essentially unique)
morphism $Y_i \to Z_0$. Since this is a morphism over $X_0$
it induces a morphism into $Z_i = Z_0 \times_{X_0} X_i$ as desired.
\end{proof}

\end{document}
